\documentclass[11pt]{amsart}
\usepackage[utf8]{inputenc}
\usepackage[T1]{fontenc}
\usepackage{lmodern}
\usepackage{amsmath,amssymb,amsthm}
\usepackage{booktabs}
\usepackage{enumitem}
\usepackage[margin=1.15in]{geometry}
\usepackage[hidelinks]{hyperref}
\usepackage{xcolor}
\usepackage{microtype}
\emergencystretch=3em

\newtheorem{theorem}{Theorem}[section]
\newtheorem{proposition}[theorem]{Proposition}
\newtheorem{lemma}[theorem]{Lemma}
\theoremstyle{definition}
\newtheorem{definition}[theorem]{Definition}
\theoremstyle{remark}
\newtheorem{remark}[theorem]{Remark}

\newcommand{\R}{\mathbb{R}}
\newcommand{\N}{\mathbb{N}}
\newcommand{\Z}{\mathbb{Z}}
\newcommand{\Q}{\mathbb{Q}}
\newcommand{\lean}[1]{\texttt{#1}}
\newcommand{\dd}{\,\mathrm{d}}

\title[A kernel-checked bound for the planar centred maximal constant]{Charge kernels and a kernel-checked
bound $c_2 \le 3.7265$ for the centred maximal operator over squares}
\author{Fernando}
\date{November 2026}

\begin{document}

\begin{abstract}
Let $c_2$ be the weak type $(1,1)$ constant of the centred Hardy--Littlewood maximal operator
over axis-parallel squares in the plane. Using Lin's comparison principle for the Cauchy
generator, we introduce \emph{charge kernels}: positive parts of $a/(|u|+|v|) - \sum_j b_j /
(|u-p_j|+|v-q_j|) - c$, whose admissibility follows from the fact that $1/(|u|+|v|)$ is a
fundamental solution of the generator, together with Kato's inequality. An $820$-charge kernel
found by linear programming, with an exact quadtree certificate of $631{,}342$ cells, gives
\[ c_2 \le 3.7265 . \]
The whole argument, analytic and computational, is formalized in Lean~4 with Mathlib, and is
checked by Lean's kernel alone: the proof uses only the standard axioms and no
\lean{native\_decide}. To make the certificate feasible for the kernel we evaluate it with a
natural-number checker that we prove computes \emph{exactly} the values of the rational one; the whole
build takes about $20$ CPU-hours. Better upper bounds are known, obtained with fractional generators:
$c_2 < 3.616$ (Lin) and $c_2 \le 3.59123$ (Tomar), the latter also proved in Lean with kernel-only
checks. The interest of the present bound lies in the different construction and in the much smaller
kernel-checked certificate.
\end{abstract}

\maketitle

\section{Introduction}

For $f \in L^1(\R^d)$ let
\[ Mf(x) = \sup_{r>0} \frac{1}{|Q(x,r)|} \int_{Q(x,r)} |f|, \qquad Q(x,r) = \prod_{i=1}^d [x_i - r, x_i + r], \]
be the centred maximal function over cubes, and let $c_d$ be the least constant with
\[ \alpha\, \bigl|\{ x : Mf(x) > \alpha \}\bigr| \le c_d \, \|f\|_1 \qquad (f \in L^1(\R^d),\ \alpha > 0). \]
Melas determined $c_1 = (11+\sqrt{61})/12 \approx 1.5675$ \cite{Melas}; no best constant is known in
higher dimension, and Aldaz proved $c_d \to \infty$ \cite{Aldaz11}. In the plane, a covering
argument gives $c_2 \le 4$, and Aldaz gave $c_2 \ge 1.6211$ \cite{Aldaz00}. In 2026 Lin proved
$1.6855 \le c_2 \le 3.879$, both formalized in Lean~4 \cite{LinRepo}; these are the values recorded in
the list of optimization constants maintained by Tao and others \cite{TaoList}. Lin also obtained
$c_2 < 3.616$ with a fractional generator \cite{LinArticle}; a Lean version of that proof existed in an
earlier state of his repository \cite{LinPB} and has since been removed from its scope. Tomar obtained
$1.7027 \le c_2 \le 3.59123$ \cite{Tomar}, with a fractional generator of order $47/40$ and a spline
kernel: the upper bound is checked with interval arithmetic and, building on Lin's formalization, also
proved in Lean with kernel-only checks (a build of several hundred CPU-hours), as is the lower bound.

This note contributes the following.
\begin{enumerate}[leftmargin=2em]
\item A family of comparison kernels for Lin's method with the Cauchy generator, \emph{charge
  kernels} (Section~\ref{sec:charges}), whose admissibility is automatic: it reduces to two
  explicit inequalities for an elementary function.
\item The bound $c_2 \le 7453/2000 = 3.7265$ (Theorem~\ref{thm:main}).
\item A complete formalization in Lean~4 \cite{Lean4} with Mathlib \cite{Mathlib}, building on
  Lin's \cite{LinRepo}, which is checked by the kernel alone. Every number in the certificate is
  recomputed by the kernel; nothing relies on the compiler (\lean{native\_decide}).
\item A technique that makes such certificates cheap for the kernel: an \emph{exact}
  natural-number refinement of a rational checker (Section~\ref{sec:nat}).
\end{enumerate}
The bound $3.7265$ is weaker than those of \cite{LinArticle,Tomar}, and Tomar's is also verified
in Lean, so this note does not improve the best known or the best formally verified bound. For the
Cauchy generator the method itself cannot reach them: a direct discretization of the optimization
problem suggests that the best kernel in this setting gives about $3.64$ (Section~\ref{sec:limits}).
What we offer is a different construction with automatic admissibility, and a kernel-checked certificate
of about $20$ CPU-hours, roughly an order of magnitude below the several hundred reported for the
Lean proof of \cite{Tomar}, made possible by the exact natural-number refinement.

\section{The comparison principle}

Write points of $\R^2$ as $z = (u,v)$, let $r(z) = |u| + |v|$, and let
$D = \{ r \le 1 \}$ be the closed unit diamond. The linear map $(u,v) \mapsto (u+v, u-v)$ maps
diamonds onto squares, so $c_2$ is also the weak type constant of the centred maximal operator over
dilates of $D$. Let $A$ be the sum of the two one-dimensional Cauchy jump generators,
\[ A\varphi(z) = \sum_{j=1}^{2} \int_{\R} \bigl( \varphi(z + t e_j) + \varphi(z - t e_j) - 2\varphi(z) \bigr)\, \frac{\dd t}{t^2}, \]
which is $-2\pi(|D_u| + |D_v|)$ in Fourier terms. The following form of Lin's comparison principle
is the one we formalized (\lean{weakTypeConstant\_two\_le\_half\_integral\_of\_pos}); Lin proved it
for one particular kernel, and the proof uses only the stated properties.

\begin{theorem}[comparison principle]\label{thm:comparison}
Let $K : \R^2 \to [0,\infty)$ be even, integrable, compactly supported, with $K \ge 1$ on $D$, and
suppose that for some $a \ge 0$
\[ \int K\, A\varphi \;\ge\; -a\,\varphi(0) \qquad \text{for every test function } \varphi \ge 0, \]
that is, $AK + a\delta_0 \ge 0$ as a distribution. Then $c_2 \le \frac12 \int K$.
\end{theorem}

The proof solves an obstacle problem $Au = \sigma - f$ with $0 \le \sigma \le \kappa$ and
$\sigma = \kappa$ on $\{u > 0\}$, so that $|\{u>0\}| \le \|f\|_1/\kappa$; off this set, pairing with
the dilates of $K$ and using $AK + a\delta_0 \ge 0$ and $u \ge 0$ bounds the averages of $f$ over
diamonds. Lin's argument assumed an exact identity $AK = \beta - a\delta_0$ with a finite positive
measure $\beta$. We only assume the inequality, which is what Kato's inequality delivers below; the
transfer step goes through with an inequality in place of the identity (\lean{Transfer.lean}).

\section{Charge kernels}\label{sec:charges}

Let $G(z) = 1/r(z)$. Since $\widehat{G}(\xi,\eta)$ is a multiple of $1/(|\xi|+|\eta|) =
\int_0^\infty e^{-t|\xi|} e^{-t|\eta|} \dd t$, $G$ is $2\pi$ times the potential of the semigroup
generated by $A$: in Lin's formalization, $\frac{1}{2\pi}\int G\, A\varphi = -\varphi(0)$ for every
test function, that is, $AG = -2\pi \delta_0$.

\begin{definition}
Let $a \ge 0$, $c > 0$, $m > 0$, and let $b_1, \dots, b_n \ge 0$ be charges at points $p_1, \dots, p_n$.
Put
\[ \psi(z) = \frac{a}{r(z)} - \sum_{j=1}^n \frac{b_j}{r(z - p_j)} - c, \qquad K = \frac{\psi_+}{m} . \]
\end{definition}

Formally $A\psi = -2\pi a\delta_0 + 2\pi \sum_j b_j \delta_{p_j}$: the charges contribute positive
point masses, and the charges sit where $\psi = -\infty$, so they do not meet $\{\psi > 0\}$.
Kato's inequality $A\psi_+ \ge \mathbf 1_{\{\psi>0\}} A\psi$ then suggests
$AK \ge -(2\pi a/m)\,\delta_0$. We prove the needed statement directly with the semigroup.

\begin{proposition}[Kato's inequality for charge kernels]\label{prop:kato}
For every test function $\varphi \ge 0$,
$\int K\, A\varphi \ge -\frac{4\pi^2 a}{m}\, \varphi(0)$.
\end{proposition}

\begin{proof}[Sketch]
Let $P_h F = p_h * F$ be the Cauchy semigroup, with $p_h(z)$ the product of two one-dimensional
Cauchy densities of scale $h$. The semigroup property $p_h * p_s = p_{h+s}$ (via the Fourier
transform) shows that the potential is excessive:
$P_h G = G - \int_0^h p_t \dd t \le G$ away from the origin. At a point where $\psi > 0$, since
$K \ge \psi/m$ and $P_h$ preserves order,
$P_h K - K \ge (P_h \psi - \psi)/m \ge -(a/m) \int_0^h p_t \dd t$ (each charge term only helps); at a
point where $\psi \le 0$, $P_h K - K \ge 0$. Integrating against $\varphi$, dividing by $h$ and using
Duhamel's formula $P_h\varphi - \varphi = \int_0^h \frac{1}{2\pi} P_s(A\varphi)\dd s$ gives the claim as
$h \to 0^+$.
\end{proof}

The remaining hypotheses of Theorem~\ref{thm:comparison} are elementary: $K \ge 0$; $K$ is even if
the configuration is symmetric under $z \mapsto -z$; $K$ has compact support, since $\psi \le a/r - c
\le 0$ for $r \ge a/c$; and $K \ge 1$ on $D$ as soon as $\psi \ge m$ there. Hence:

\begin{theorem}\label{thm:charge}
If the charge configuration is symmetric under $z \mapsto -z$, $\psi \ge m$ on $D \setminus\{0\}$, and
$\int \psi_+ \le 2mB$, then $c_2 \le B$.
\end{theorem}

The charge-free kernel ($n=0$) is Lin's tail $(R/r - 1)_+$ up to normalization; with $R = 2$ it
gives $c_2 \le 4$. Adding charges removes mass where $\psi$ is positive but large mass is not needed.

\section{Choosing the charges}\label{sec:lp}

For fixed charge positions the problem ``minimize $\int\psi_+$ subject to $\psi \ge m$ on $D$'' is
convex in $(a, c, b)$. We discretize it into a linear program (solved with HiGHS through SciPy):
the objective is written in diamond-polar coordinates (below), in which $r\psi_+ =
\max(a - r(c + S), 0)$ is bounded, with $S$ the charge sum; positions are placed on a grid outside $D$
in orbits of the symmetry group of the square. Two observations improved the first attempts:
(i) charges far from the origin (out to $r \approx 6$) with a small constant $c$, so that
$\psi_+$ has a thin tail out to $r = a/c \approx 40$; and (ii) a normalization $m$ slightly above $1$.

The kernel used below has $820$ charges: a grid of spacing $0.1$ for $1 \le r < 2.6$ and $0.5$ for
$2.6 \le r < 6$. Its data are rounded to dyadic rationals: positions in $2^{-12}\Z$, strengths in
$2^{-24}\Z$, and
\[ a = \tfrac{31624813}{2^{24}} \approx 1.88498, \qquad c = \tfrac{197655}{2^{22}} \approx 0.047124,
\qquad m = \tfrac{1003}{1000} . \]
The configuration is invariant under both reflections $u \mapsto -u$ and $v \mapsto -v$.

\section{Exact certificates}\label{sec:cert}

\subsection*{Coordinates and symmetry}
In the quadrant with signs $(\sigma_u, \sigma_v)$ the map $(r,s) \mapsto (\sigma_u rs,\ \sigma_v r(1-s))$,
$r > 0$, $0 < s < 1$, has Jacobian $r$ and satisfies $r(z) = r$, so
\[ \int \psi_+ = \sum_{\sigma_u, \sigma_v = \pm1} \int_0^\infty \!\! \int_0^1 \max\bigl(a - r(c + S),\ 0\bigr) \dd s \dd r, \]
where $S$ is the charge sum at the image point. By the two reflection symmetries (checked on the
charge list) the four quadrants contribute equally, and $\psi_+ = 0$ for $r \ge R := 64 > a/c$.

\subsection*{Leaf bounds}
A quadtree subdivides $[0, R] \times [0, 1]$ into cells $[r_0, r_1] \times [s_0, s_1]$ with dyadic corners.
The image of a cell lies in a box $[u_0,u_1] \times [v_0,v_1]$, and for each charge the $\ell^1$
distance to points of the box lies in $[d_{\min}, d_{\max}]$. Reciprocals are rounded outward to
$2^{-30}\Z$: $S \ge S_{\min} = \sum_j \lfloor b_j/d_{\max,j} \rfloor_{30}$ and
$S \le S_{\max} = \sum_j \lceil b_j/d_{\min,j} \rceil_{30}$. Each leaf gets the smaller of two bounds for
$\iint \max(a - r(c+S), 0)$:
\begin{itemize}[leftmargin=2em]
\item[$J_1$:] with $k = c + S_{\min}$, the integral of $\max(a - rk, 0)$, bounded by $0$, by $a - rk$, or by
  the chord through $(r_0, a - r_0k)$ and $(r_1, 0)$, according to the signs of $a - r_0 k$ and $a - r_1 k$;
\item[$J_2$:] if every charge is at positive distance, no charge coordinate straddles the box, and
  $a - r_1(c + S_{\max}) \ge 0$ (so the integrand is $a - r(c+S)$ on the cell), a tangent bound: on the box
  each distance is an affine function $t$ of $(u,v)$, and $b/t \ge \alpha - \beta t$ for all $t>0$ whenever
  $\alpha^2 \le 4b\beta$, $\beta \ge 0$. With rounded $\alpha \approx 2b/t_0$, $\beta \approx b/t_0^2$ at the
  centre, this gives a polynomial upper bound $e_0 + e_1 r + e_2 r^2 + e_3 r^2 s$, integrated exactly.
\end{itemize}
A second quadtree on $[0,1]^2$ certifies $\psi \ge m$ on $D$, through $a \ge r_1(c + m + S_{\max})$ on each
cell. The soundness of both checks is proved in Lean (\lean{v2\_sound}, \lean{v1\_sound}), together with
the reduction to one quadrant (\lean{c2\_le\_of\_cert\_sym}).

\subsection*{The certificate}
The integral tree has $631{,}342$ leaves (depth $14$), chosen by a floating-point refinement
that splits cells where the gap between the bound and the estimate is largest; in a random sample of
$300$ leaves, about a quarter use the tangent bound. The diamond tree has $1{,}540$ leaves. The exact value of the integral certificate for
one quadrant is $V/(6 \cdot 2^{105})$ with
\[ V = 454850744367578333886108605087744, \qquad \frac{V}{6\cdot 2^{105}} = 1.86882276\ldots, \]
so $\int\psi_+ \le 4V/(6\cdot2^{105})$, and Theorem~\ref{thm:charge} gives
\[ c_2 \;\le\; \frac{4V}{2m \cdot 6 \cdot 2^{105}} = 3.7264661\ldots \;\le\; 3.7265 . \]

\begin{theorem}\label{thm:main}
$c_2 \le 7453/2000 = 3.7265$. In Lean this is the theorem
\[ \texttt{ChargeKernel.Pt900.weakTypeConstant\_two\_le} , \]
stated with Lin's definition \lean{CenteredMaximal.weakTypeConstant 2} and proved with the axioms
\lean{propext}, \lean{Classical.choice}, \lean{Quot.sound} only.
\end{theorem}

\section{Checking the certificate in the kernel}\label{sec:nat}

The checker is a Lean function on rationals. Lean's kernel can evaluate it (\lean{decide +kernel}),
but each rational operation normalizes with a gcd and is slow in the kernel: a cell with $820$ charges
costs seconds, and the full certificate would take hundreds of CPU-hours. The usual shortcut,
\lean{native\_decide}, runs compiled code and adds the compiler to the trusted base. We avoid both.

\subsection*{An exact natural-number refinement}
All data are dyadic, so every quantity the checker computes is an integer at a fixed scale. With
$r = R_n/2^8$ and $s = S_n/2^{14}$ for the integral tree, box coordinates and distances are integers
times $2^{-22}$ (positions are stored with an offset $2^{32}$ to stay in $\N$), the rounded
reciprocals and $S_{\min}, S_{\max}$ are integers times $2^{-30}$, and both leaf bounds are integers
times $1/(6 \cdot 2^{105})$. The file \lean{CertNat.lean} re-implements the checker in these units: the
loops over charges use only operations on $\N$, which the kernel evaluates on literals with GMP,
and the few per-cell polynomial terms use $\Z$. It does not import Mathlib.
The file \lean{CertNatSound.lean} proves, for decoded data,
\[ \texttt{nChk} = \mathrm{some}\, v \;\Longrightarrow\; \texttt{v2chk} = \mathrm{some}\bigl(v/(6\cdot 2^{105})\bigr), \qquad
\texttt{nV1} = \mathrm{true} \;\Longrightarrow\; \texttt{v1ok} = \mathrm{true}, \]
covering every rounding step, both leaf bounds and the tree recursion. The soundness theorem of the
rational checker is then used unchanged. In practice the natural-number checker is about two orders
of magnitude faster in the kernel.

\subsection*{Sharding}
The integral tree is cut into subtrees of at most $32$ leaves, each checked by its own
\lean{decide +kernel} theorem with the exact integer value of the subtree; a lemma
\lean{nChk\_node} assembles four children into their parent. The $631{,}342$ leaves give $310$
modules. Their kernel checks take about $20$ CPU-hours in total (about $3.5$ hours with $8$ modules
in parallel, $1$--$2$ GB of memory each); the decoding of the charge list and the symmetry checks take
a few seconds.

\section{Trust and independent checks}

The trusted base of Theorem~\ref{thm:main} is Lean's kernel, the three standard axioms, and the
definitions in the statement. We checked that Lin's \lean{weakTypeConstant 2} is the standard
constant: \lean{maximalFunction} averages over \lean{closedBall x r} in $\R^2$ with the sup norm,
which is the closed square of side $2r$, and \lean{weakTypeConstant} is the infimum of the weak type
bounds over all integrable $f$ and all levels. The new files contain no \lean{sorry}, \lean{axiom},
\lean{implemented\_by}, \lean{extern}, \lean{unsafe}, \lean{partial} or \lean{native\_decide}.
As further checks:
\begin{itemize}[leftmargin=2em]
\item an independent transliteration of the rational checker in Python, in exact arithmetic
  (\texttt{checker\_rational.py}), reproduces the certificate value $V/(6\cdot2^{105})$ exactly;
\item the generator of the Lean shards recomputes every leaf value in Python and asserts it is
  nonnegative before the kernel checks it again;
\item an earlier version of the formalization (all modules except the $310$ shards) was replayed
  with Lean's standalone checker \lean{leanchecker}.
\end{itemize}
Mathlib was used from its prebuilt cache and was not replayed from source, and the shard modules have
not been replayed with an independent kernel implementation; both are natural further checks.

\section{Limits of the method and other remarks}\label{sec:limits}

\subsection*{The Cauchy generator}
A direct linear program over all kernels on a grid, with the operator discretized, gives
$\frac12\int K = 3.6744$, $3.6461$, $3.6416$, $3.6403$ for mesh sizes $1/10$, $1/20$, $1/30$, $1/40$
(\texttt{kernel\_lp\_sym.py}). This suggests that the comparison principle with the Cauchy generator
cannot give much less than $3.64$. The optimal grid kernel is not a function of $r$: it extends to
about $1.43$ along the axes and barely beyond $D$ along the diagonals. The better bounds of
\cite{LinArticle,Tomar} use generators of order $\alpha \in (1,2)$.

\subsection*{Gap between estimate and certificate}
The floating-point estimate for our $820$-charge kernel is $3.697$; the certificate loses $0.03$, mostly
(about three quarters) in cells crossing the zero set of $\psi$, where only $J_1$ applies. Replacing
point charges by charges spread uniformly over small squares, whose potentials have closed forms
from the antiderivative $(x+y)\log(x+y)$, gives floating-point estimates down to $3.671$; formalizing
these would need Kato's inequality for an absolutely continuous charge distribution.

\subsection*{Lower bounds}
We reproduced Lin's lower bound $1.6855$ exactly but did not improve it: searches over periodic
configurations with up to nine atoms per cell, product families and two-scale constructions stayed
at or below it. Tomar's $1.7027$ \cite{Tomar} uses a periodic measure with $50$ atoms per cell.

\section{Reproducibility}

The repository accompanying this note contains the Lean development (which depends on Lin's
repository at a pinned commit), the kernel and certificate data, the scripts that produced them, the
two Python reference checkers, and the generator of the Lean shards. A clean build is
\texttt{cd lean \&\& ./build.sh}, which fetches Mathlib's cache, builds the shards as a
memory-bounded work queue, builds the main theorem and prints its axioms.

\subsection*{Acknowledgements}
This work builds directly on Yongxi Lin's Lean formalization \cite{LinRepo}: the definition of
$c_2$, the Cauchy generator and its potential, and the comparison argument. The research,
the Lean development and this text were produced with the assistance of Claude (Anthropic), an AI
system, under the author's direction.

\begin{thebibliography}{99}
\bibitem{Aldaz00} J. M. Aldaz, \emph{A remark on the centered $n$-dimensional Hardy--Littlewood maximal
  function}, Czechoslovak Math. J. \textbf{50} (2000), 103--112.
\bibitem{Aldaz11} J. M. Aldaz, \emph{The weak type $(1,1)$ bounds for the maximal function associated to
  cubes grow to infinity with the dimension}, Ann. of Math. \textbf{173} (2011), 1013--1023.
\bibitem{Lean4} L. de Moura and S. Ullrich, \emph{The Lean 4 theorem prover and programming language},
  CADE-28, Lecture Notes in Comput. Sci. \textbf{12699} (2021), 625--635.
\bibitem{LinRepo} Y. Lin, \emph{Centred maximal constants for squares and Euclidean balls}, Lean~4
  formalization, \url{https://github.com/CoolRmal/centered-maximal-constant} (2026).
\bibitem{LinArticle} Y. Lin, \emph{Two partial balayage principles} (2026),
  \url{https://coolrmal.github.io/articles/two-partial-balayage-principles/}.
\bibitem{LinPB} Y. Lin, \emph{Two partial balayage principles in Lean},
  \url{https://github.com/CoolRmal/partial-balayage-lean} (2026).
\bibitem{Mathlib} The mathlib Community, \emph{The Lean mathematical library}, CPP 2020, 367--381.
\bibitem{Melas} A. D. Melas, \emph{The best constant for the centered Hardy--Littlewood maximal
  inequality}, Ann. of Math. \textbf{157} (2003), 647--688.
\bibitem{TaoList} T. Tao et al., \emph{Optimization constants in mathematics}, entry 47a,
  \url{https://teorth.github.io/optimizationproblems/constants/47a.html}.
\bibitem{Tomar} T. S. Tomar, \emph{Bounds for the weak type $(1,1)$ constant of the centred maximal operator
  over squares}, Zenodo, \href{https://doi.org/10.5281/zenodo.23543279}{doi:10.5281/zenodo.23543279};
  certificates, checkers and Lean proofs at \url{https://github.com/tejasvi-releases/square-maximal-weak-type-v3}.
\end{thebibliography}

\end{document}
